\documentclass[11pt,fontset=fandol]{ctexart}
\usepackage[a4paper,margin=23mm]{geometry}
\usepackage{amsmath,amssymb,mathtools,bm,mathrsfs,esint}
\usepackage{graphicx,longtable,booktabs,array,calc}
\usepackage[normalem]{ulem}
\usepackage{xcolor,hyperref}
\hypersetup{hidelinks}
\setcounter{secnumdepth}{-1}
\newcounter{none}
\providecommand{\tightlist}{\setlength{\itemsep}{0pt}\setlength{\parskip}{0pt}}
\providecommand{\passthrough}[1]{#1}
\usepackage{amssymb,mathrsfs,esint,bm,newunicodechar}
\allowdisplaybreaks
\setlength{\emergencystretch}{3em}
\providecommand{\llangle}{\langle\!\langle}
\providecommand{\rrangle}{\rangle\!\rangle}
\newunicodechar{∞}{\ensuremath{\infty}}
\newunicodechar{∈}{\ensuremath{\in}}
\newunicodechar{∑}{\ensuremath{\sum}}
\newunicodechar{∫}{\ensuremath{\int}}
\newunicodechar{≤}{\ensuremath{\leq}}
\newunicodechar{≥}{\ensuremath{\geq}}
\newunicodechar{≠}{\ensuremath{\neq}}
\newunicodechar{→}{\ensuremath{\to}}
\newunicodechar{⇒}{\ensuremath{\Rightarrow}}
\newunicodechar{α}{\ensuremath{\alpha}}
\newunicodechar{β}{\ensuremath{\beta}}
\newunicodechar{γ}{\ensuremath{\gamma}}
\newunicodechar{λ}{\ensuremath{\lambda}}
\newunicodechar{φ}{\ensuremath{\varphi}}
\newunicodechar{ε}{\ensuremath{\varepsilon}}
\newunicodechar{σ}{\ensuremath{\sigma}}
\newunicodechar{θ}{\ensuremath{\theta}}
\newunicodechar{δ}{\ensuremath{\delta}}
\newunicodechar{η}{\ensuremath{\eta}}
\newunicodechar{π}{\ensuremath{\pi}}
\newunicodechar{ξ}{\ensuremath{\xi}}
\newunicodechar{Σ}{\ensuremath{\Sigma}}
\newunicodechar{Δ}{\ensuremath{\Delta}}
\newunicodechar{∂}{\ensuremath{\partial}}
\newunicodechar{∀}{\ensuremath{\forall}}
\newunicodechar{∃}{\ensuremath{\exists}}
\newunicodechar{√}{\ensuremath{\surd}}
\newunicodechar{−}{\ensuremath{-}}
\newunicodechar{μ}{\ensuremath{\mu}}
\newunicodechar{ρ}{\ensuremath{\rho}}
\newunicodechar{τ}{\ensuremath{\tau}}
\newunicodechar{ω}{\ensuremath{\omega}}
\newunicodechar{Ω}{\ensuremath{\Omega}}
\newunicodechar{χ}{\ensuremath{\chi}}
\newunicodechar{ψ}{\ensuremath{\psi}}
\newunicodechar{Γ}{\ensuremath{\Gamma}}
\newunicodechar{Φ}{\ensuremath{\Phi}}
\newunicodechar{Λ}{\ensuremath{\Lambda}}
\newunicodechar{ℝ}{\ensuremath{\mathbb{R}}}
\newunicodechar{ℤ}{\ensuremath{\mathbb{Z}}}
\newunicodechar{ℕ}{\ensuremath{\mathbb{N}}}
\newunicodechar{△}{\ensuremath{\triangle}}
\newunicodechar{⊗}{\ensuremath{\otimes}}
\begin{document}
\section{2023年高等代数参考答案}\label{ux5e74ux9ad8ux7b49ux4ee3ux6570ux53c2ux8003ux7b54ux6848}

来源：\texttt{HDU/601数学分析/2023/2023年杭州电子科技大学数分高代真题及答案(双科合卷).pdf}。原
PDF 页序：5, 6, 7, 8, 9, 10, 11, 12, 13。

\section{杭州电子科技大学2023年高等代数真题解析}\label{ux676dux5ddeux7535ux5b50ux79d1ux6280ux5927ux5b662023ux5e74ux9ad8ux7b49ux4ee3ux6570ux771fux9898ux89e3ux6790}

\subsection{一、计算行列式}\label{ux4e00ux8ba1ux7b97ux884cux5217ux5f0f}

先由下到上作相邻行相减，再对得到的后四行继续由下到上作相邻行相减，并消去第二行的首项，可将原解析缺损的行列式运算恢复为
\[
\begin{aligned}
D&=\begin{vmatrix}
1&2&3&4&5\\1&1&1&1&-4\\1&1&1&-4&1\\1&1&-4&1&1\\1&-4&1&1&1
\end{vmatrix}\\
&=\begin{vmatrix}
1&2&3&4&5\\0&-1&-2&-3&-9\\0&0&0&-5&5\\0&0&-5&5&0\\0&-5&5&0&0
\end{vmatrix}\\
&=\begin{vmatrix}-1&-2&-3&-9\\0&0&-5&5\\0&-5&5&0\\-5&5&0&0\end{vmatrix}.
\end{aligned}
\] 依次令第二列加第一列、第三列加第二列、第四列加第三列，得 \[
\begin{aligned}
D&=\begin{vmatrix}-1&-3&-3&-9\\0&0&-5&5\\0&-5&5&0\\-5&0&0&0\end{vmatrix}\\
&=\begin{vmatrix}-1&-3&-6&-9\\0&0&-5&5\\0&-5&0&0\\-5&0&0&0\end{vmatrix}\\
&=\begin{vmatrix}-1&-3&-6&-15\\0&0&-5&0\\0&-5&0&0\\-5&0&0&0\end{vmatrix}\\
&=(-1)^{4\cdot3/2}(-15)(-5)^3=1875.
\end{aligned}
\] 以上补回了OCR漏失的零元与等号，并删除由水印造成的虚假重复行。

\subsection{二、矩阵方程}\label{ux4e8cux77e9ux9635ux65b9ux7a0b}

（1）行变换为 \[
A\longrightarrow\begin{pmatrix}1&1&0&-3\\0&-6&6&15\end{pmatrix}
\longrightarrow\begin{pmatrix}1&0&1&-1/2\\0&1&-1&-5/2\end{pmatrix}.
\] 故 \(Ax=0\) 的通解为 \[
\begin{aligned}
x&=\begin{pmatrix}-x_3+x_4/2\\x_3+5x_4/2\\x_3\\x_4\end{pmatrix}\\
&=x_3\begin{pmatrix}-1\\1\\1\\0\end{pmatrix}
+\frac{x_4}{2}\begin{pmatrix}1\\5\\0\\2\end{pmatrix}.
\end{aligned}
\] 因此取 \[
B=\begin{pmatrix}-1&1\\1&5\\1&0\\0&2\end{pmatrix},
\] 即可满足 \(AB=0\)，且两列线性无关，\(r(B)=2\)。

（2）增广矩阵必须有六列，其行变换为 \[
\begin{aligned}
(A,E)&=\left(\begin{array}{rrrr|rr}1&1&0&-3&1&0\\4&-2&6&3&0&1\end{array}\right)\\
&\longrightarrow\left(\begin{array}{rrrr|rr}1&1&0&-3&1&0\\0&-6&6&15&-4&1\end{array}\right)\\
&\longrightarrow\left(\begin{array}{rrrr|rr}1&1&0&-3&1&0\\0&1&-1&-5/2&2/3&-1/6\end{array}\right)\\
&\longrightarrow\left(\begin{array}{rrrr|rr}1&0&1&-1/2&1/3&1/6\\0&1&-1&-5/2&2/3&-1/6\end{array}\right).
\end{aligned}
\] 所以 \(AX=E\) 的通解为 \[
X=\begin{pmatrix}
-k_1+l_1/2+1/3&-k_2+l_2/2+1/6\\
k_1+5l_1/2+2/3&k_2+5l_2/2-1/6\\
k_1&k_2\\l_1&l_2
\end{pmatrix}.
\] 其中 \(k_1,k_2,l_1,l_2\) 为任意数。特别可取 \[
X=\begin{pmatrix}1/3&1/6\\2/3&-1/6\\0&0\\0&0\end{pmatrix},\qquad AX=E_2.
\]

\subsection{三、线性方程组}\label{ux4e09ux7ebfux6027ux65b9ux7a0bux7ec4}

（1）系数矩阵为 \[
M=\begin{pmatrix}\lambda&1&1&1\\1&\lambda&1&1\\1&1&\lambda&1\\1&1&1&\lambda\end{pmatrix}
=(\lambda-1)E+\begin{pmatrix}1&1&1&1\\1&1&1&1\\1&1&1&1\\1&1&1&1\end{pmatrix}.
\] 特征值为 \(\lambda-1\)（三重）、\(\lambda+3\)，故 \(\lambda\ne1,-3\)
时有唯一解。设 \(S=x_1+x_2+x_3+x_4\)，将四个方程相加，得到 \[
(\lambda+3)S=2\lambda+2,\qquad S=\frac{2\lambda+2}{\lambda+3}.
\] 分别从原方程左侧减去 \(S\)，有 \[
\begin{aligned}
x_1=x_2&=\frac{1-S}{\lambda-1}=-\frac1{\lambda+3},\\
x_3=x_4&=\frac{\lambda-S}{\lambda-1}=\frac{\lambda+2}{\lambda+3}.
\end{aligned}
\] 唯一解为 \[
x=\left(-\frac1{\lambda+3},-\frac1{\lambda+3},\frac{\lambda+2}{\lambda+3},\frac{\lambda+2}{\lambda+3}\right)^T.
\]

（2）\(\lambda=1\) 时，方程组等价于 \(x_1+x_2+x_3+x_4=1\)，通解为 \[
\begin{aligned}
x&=\begin{pmatrix}x_1\\x_2\\x_3\\1-x_1-x_2-x_3\end{pmatrix}\\
&=x_1\begin{pmatrix}1\\0\\0\\-1\end{pmatrix}
+x_2\begin{pmatrix}0\\1\\0\\-1\end{pmatrix}
+x_3\begin{pmatrix}0\\0\\1\\-1\end{pmatrix}
+\begin{pmatrix}0\\0\\0\\1\end{pmatrix}.
\end{aligned}
\] \(\lambda=-3\) 时四个方程左侧之和为零，右侧之和为
\(2\lambda+2=-4\ne0\)，故无解。

\subsection{四、二次型}\label{ux56dbux4e8cux6b21ux578b}

二次型的矩阵为 \[
H=\begin{pmatrix}2&2&-2\\2&5&-4\\-2&-4&5\end{pmatrix}
=E+\begin{pmatrix}1&2&-2\\2&4&-4\\-2&-4&4\end{pmatrix}
=E+\begin{pmatrix}1\\2\\-2\end{pmatrix}(1,2,-2).
\] 由秩一矩阵性质，特征值为 \(1,1,1+(1,2,-2)(1,2,-2)^T=10\)。特征值
\(1\) 对应方程 \((1,2,-2)x=0\) 的非零解，取两条正交特征向量 \[
v_1=(0,1,1)^T,\qquad v_2=(4,-1,1)^T.
\] 特征值 \(10\) 对应 \(v_3=(1,2,-2)^T\)。单位化后取 \[
T=\begin{pmatrix}
0&4/(3\sqrt2)&1/3\\
1/\sqrt2&-1/(3\sqrt2)&2/3\\
1/\sqrt2&1/(3\sqrt2)&-2/3
\end{pmatrix}.
\] 则 \(T\) 正交，\(T^THT=\operatorname{diag}(1,1,10)\)，标准形为
\(f=y_1^2+y_2^2+10y_3^2\)。

\subsection{五、线性变换}\label{ux4e94ux7ebfux6027ux53d8ux6362}

（1）以新基的旧基坐标为列，过渡矩阵为 \[
P=\begin{pmatrix}1&0&0\\-2&3&0\\0&-1&1\end{pmatrix},\qquad
P^{-1}=\frac13\begin{pmatrix}3&0&0\\2&1&0\\2&1&3\end{pmatrix}.
\] 新基下的矩阵为 \[
\begin{aligned}
B&=P^{-1}AP\\
&=\frac13\begin{pmatrix}3&0&0\\2&1&0\\2&1&3\end{pmatrix}
\begin{pmatrix}1&-2&2\\-5&5&1\\-3&1&5\end{pmatrix}\\
&=\frac13\begin{pmatrix}3&-6&6\\-3&1&5\\-12&4&20\end{pmatrix}.
\end{aligned}
\]

（2） \[
A\longrightarrow\begin{pmatrix}1&0&2\\0&2&3\\0&2&3\end{pmatrix}
\longrightarrow\begin{pmatrix}1&0&2\\0&1&3/2\\0&0&0\end{pmatrix}.
\] 故 \(A\) 的前两列为极大线性无关组，值域为 \[
\mathcal AV=\operatorname{span}(\varepsilon_1-\varepsilon_2+\varepsilon_3,\ 2\varepsilon_2+2\varepsilon_3).
\] 解 \(Ax=0\) 得 \[
x=\begin{pmatrix}-2x_3\\-3x_3/2\\x_3\end{pmatrix}
=\frac{x_3}{2}\begin{pmatrix}-4\\-3\\2\end{pmatrix}.
\] 因此 \[
\mathcal A^{-1}(\theta)=\ker\mathcal A
=\operatorname{span}(-4\varepsilon_1-3\varepsilon_2+2\varepsilon_3).
\]

\subsection{六、矩阵求逆}\label{ux516dux77e9ux9635ux6c42ux9006}

已知 \(D=E+BCA\)、\(C=(E-AB)^{-1}\)。

（1） \[
\begin{aligned}
E&=E-BA+BA\\
&=E-BA+B(E-AB)(E-AB)^{-1}A\\
&=E-BA+B(E-AB)^{-1}A-BAB(E-AB)^{-1}A\\
&=[E+B(E-AB)^{-1}A]-BA[E+B(E-AB)^{-1}A]\\
&=(E-BA)[E+B(E-AB)^{-1}A].
\end{aligned}
\] 方阵有右逆即为可逆，所以 \[
(E-BA)^{-1}=E+B(E-AB)^{-1}A=E+BCA=D.
\]

（2） \[
\begin{aligned}
AD&=A+AB(E-AB)^{-1}A\\
&=A+(AB-E+E)(E-AB)^{-1}A\\
&=A+[-E+(E-AB)^{-1}]A=(E-AB)^{-1}A.
\end{aligned}
\] 因此 \[
ADB=(E-AB)^{-1}AB=(E-AB)^{-1}(AB-E+E)=-E+C,
\] 即 \(C=E+ADB\)。

\subsection{七、公共特征向量}\label{ux4e03ux516cux5171ux7279ux5f81ux5411ux91cf}

（1）任取正整数 \(k\)，利用迹的循环性质，有 \[
\begin{aligned}
\operatorname{tr}(A^k)
&=\operatorname{tr}(A^{k-1}(AB-BA))\\
&=\operatorname{tr}(A^kB)-\operatorname{tr}(A^{k-1}BA)\\
&=\operatorname{tr}(A^kB)-\operatorname{tr}(AA^{k-1}B)=0.
\end{aligned}
\] 特别对 \(k=1,\ldots,n\) 成立。由特征值的 Newton
恒等式，所有初等对称式均为零，所以特征多项式为 \(\lambda^n\)，\(A\)
的所有特征值均为零。

（2）记 \(W=\ker A\)，由（1）知 \(W\ne\{\theta\}\)。若 \(x\in W\)，则 \[
ABx=BAx+Ax=0,
\] 所以 \(Bx\in W\)，\(W\) 是 \(B\) 的不变子空间。将 \(B\)
限制在非零复向量空间 \(W\) 上，必有特征向量 \(x\ne0\)，满足
\(Bx=\lambda x\)，而 \(Ax=0\)。这就是公共特征向量。原OCR在此混淆 \(W\)
与 \(V\)，以上按所定义的特征子空间恢复。

\subsection{八、实根多项式}\label{ux516bux5b9eux6839ux591aux9879ux5f0f}

（1）\(a_i\) 为 \(f\) 的单根，故 \(f'(a_i)\ne0\)，从而 \[
[f'(a_i)]^2-f(a_i)f''(a_i)=[f'(a_i)]^2>0.
\] 对 \(x\ne a_1,\ldots,a_n\)，对乘积求导得 \[
\frac{f'(x)}{f(x)}=\sum_{i=1}^n\frac1{x-a_i}.
\] 于是 \[
\begin{aligned}
\frac{[f'(x)]^2-f(x)f''(x)}{f(x)^2}
&=-\left(\frac{f'(x)}{f(x)}\right)'\\
&=-\left(\sum_{i=1}^n\frac1{x-a_i}\right)'\\
&=\sum_{i=1}^n\frac1{(x-a_i)^2}>0.
\end{aligned}
\] 结合根处的计算，所给多项式在全体实数上为正，没有实根。

（2）乘积求导给出 \[
f'(x)=\sum_{i=1}^n\prod_{j\ne i}(x-a_j)
=\left(\sum_{i=1}^n\frac1{x-a_i}\right)f(x),\quad x\ne a_i.
\] 对 \(x\ne0,a_i\)，用有限等比数列公式， \[
\begin{aligned}
&\frac1{x-a_i}\left(\frac{a_i}{x}\right)^{k+1}
+\frac1x\sum_{r=0}^k\left(\frac{a_i}{x}\right)^r\\
&=\frac1{x-a_i}\left(\frac{a_i}{x}\right)^{k+1}
+\frac1x\frac{1-(a_i/x)^{k+1}}{1-a_i/x}\\
&=\frac1{x-a_i}\left(\frac{a_i}{x}\right)^{k+1}
+\frac{1-(a_i/x)^{k+1}}{x-a_i}
=\frac1{x-a_i}.
\end{aligned}
\] 对 \(i\) 求和后乘以 \(f(x)\)，即得题中等式。原OCR把外层
\(\sum_{i=1}^n\) 误作 \(\sum_{k=1}^n\)，与固定参数 \(k\)
冲突；这里根据逐项恒等式恢复一致下标。

\subsection{九、根子空间分解}\label{ux4e5dux6839ux5b50ux7a7aux95f4ux5206ux89e3}

先证明：若 \(f,g\in P[x]\) 互素，\(h=fg\)，\(\sigma\) 为线性变换，则 \[
\ker h(\sigma)=\ker f(\sigma)\oplus\ker g(\sigma).
\] 由互素性存在 \(u,v\in P[x]\) 使 \(uf+vg=1\)，从而 \[
u(\sigma)f(\sigma)+v(\sigma)g(\sigma)=E.
\] 任取 \(x\in\ker h(\sigma)\)， \[
x=u(\sigma)f(\sigma)x+v(\sigma)g(\sigma)x.
\] 由于同一变换的多项式彼此可交换， \[
\begin{aligned}
g(\sigma)u(\sigma)f(\sigma)x&=u(\sigma)h(\sigma)x=0,\\
f(\sigma)v(\sigma)g(\sigma)x&=v(\sigma)h(\sigma)x=0.
\end{aligned}
\] 所以第一项在 \(\ker g(\sigma)\) 中，第二项在 \(\ker f(\sigma)\)
中。反过来这两个核都包含于 \(\ker h(\sigma)\)，故二者之和等于
\(\ker h(\sigma)\)。

若 \(x\in\ker f(\sigma)\cap\ker g(\sigma)\)，则 \[
x=u(\sigma)f(\sigma)x+v(\sigma)g(\sigma)x=0,
\] 所以两核交集为零，和为直和。

回到本题，由 Cayley--Hamilton 定理， \[
f(A)=(A-\lambda_1E)^{r_1}\cdots(A-\lambda_sE)^{r_s}=0.
\] 各因子 \((\lambda-\lambda_i)^{r_i}\)
两两互素，迭代使用刚才的结论，得到 \[
\begin{aligned}
V=\ker f(A)
&=\ker(A-\lambda_1E)^{r_1}\oplus
\ker\left[\prod_{i=2}^s(A-\lambda_iE)^{r_i}\right]\\
&=\ker(A-\lambda_1E)^{r_1}\oplus\ker(A-\lambda_2E)^{r_2}
\oplus\ker\left[\prod_{i=3}^s(A-\lambda_iE)^{r_i}\right]\\
&=\bigoplus_{i=1}^s\ker(A-\lambda_iE)^{r_i}
=V_1\oplus\cdots\oplus V_s.
\end{aligned}
\] 中间分步式在相应因子存在时使用；\(s=1\) 时由 \(f(A)=0\) 直接得到
\(V=V_1\)。
\end{document}
